\documentclass[11pt]{article}

\usepackage[margin=1in]{geometry}
\usepackage{amsmath,amssymb,amsthm,mathtools}
\usepackage{booktabs}
\usepackage{array}
\usepackage{enumitem}
\usepackage[numbers,sort&compress]{natbib}
\usepackage{hyperref}
\usepackage{microtype}
\usepackage{xcolor}

\hypersetup{colorlinks=true,linkcolor=blue,citecolor=blue,urlcolor=blue}
\allowdisplaybreaks

\newtheorem{proposition}{Proposition}
\newtheorem{definition}{Definition}
\newtheorem{remark}{Remark}
\numberwithin{equation}{section}

\newcommand{\R}{\mathbb R}
\newcommand{\E}{\mathbb E}
\newcommand{\N}{\mathcal N}
\newcommand{\dd}{\,\mathrm d}
\newcommand{\tr}{\operatorname{tr}}
\newcommand{\diag}{\operatorname{diag}}
\newcommand{\chol}{\operatorname{chol}}
\newcommand{\ellhat}{\widehat\ell}
\newcommand{\sgq}{\operatorname{SGQ}}
\newcommand{\branch}{\mathfrak B}

\title{Implementation-Ready Mathematical Companion to Fixed Sparse-Grid Quadrature Filtering and Its Analytical Gradient}
\author{Chak Wong}
\date{2026-06-06}

\begin{document}
\maketitle

\begin{abstract}
This note defines a fixed sparse-grid quadrature filter (FixedSGQF) for
high-dimensional nonlinear state-space models in which exact nonlinear filtering
is unavailable but a deterministic Gaussian-surrogate likelihood is still
scientifically useful.  The method carries one Gaussian state
\(\N(m_t,P_t)\), evaluates the nonlinear moments needed for Gaussian projection
with a fixed sparse-grid cloud in standardized Gaussian coordinates, and defines
an analytical gradient of the same deterministic approximate likelihood scalar
that the value routine evaluates.  The note is implementation-ready at the
algorithm level: it states the cloud-construction rule, the one-step value
recursion, the branch-conditioned derivative recursion, the same-scalar
finite-difference validity rule, and a worked numerical oracle.  The main
limitation is explicit.  FixedSGQF is not exact nonlinear filtering and it does
not preserve general non-Gaussian posterior shape; it is a deterministic
Gaussian-surrogate lane whose usefulness depends on whether one carried Gaussian
and a fixed sparse-grid moment engine are adequate for the intended regime.
\end{abstract}

\tableofcontents

\section{Reader Orientation, Scope, and Approval Target}
\label{sec:reader-orientation}

This document is written for three readers at once:
\begin{enumerate}[leftmargin=0.3in]
  \item implementation engineers who need an auditable mathematical contract,
  \item mathematically trained first-year-PhD-level workers who must implement
  the algorithm correctly from the note, and
  \item general scientific reviewers who need to judge whether the proposal is
  coherent, bounded, and worth approval.
\end{enumerate}

The note therefore serves two roles at once.  It is a mathematical derivation of
one deterministic approximate filtering lane, and it is an implementation handoff
document for that lane.  It is not a general survey of nonlinear filtering, and
it is not a full software engineering manual.

\paragraph{Background assumed.}
The reader is assumed to know multivariate Gaussian notation, basic
state-space-model notation, Jacobian-style derivative notation, and the use of a
Cholesky factor for a positive-definite covariance matrix.  The note does not
assume prior familiarity with fixed-branch derivative contracts, sparse-grid
cloud merging, or same-scalar finite-difference validation; those are defined
here.

\paragraph{What this note provides internally.}
The note defines:
\begin{enumerate}[leftmargin=0.3in]
  \item the exact filtering target and the narrower approximate object carried by
  FixedSGQF,
  \item the fixed sparse-grid cloud construction in standardized coordinates,
  \item the one-step value recursion and the accumulated deterministic
  approximate log likelihood,
  \item the fixed-branch analytical derivative of that same scalar,
  \item the same-scalar finite-difference validity rule, and
  \item a worked numerical oracle for value and gradient checking.
\end{enumerate}

\paragraph{What this note does not claim.}
The note does not claim exact nonlinear filtering, exact nonlinear likelihood
evaluation, or faithful preservation of arbitrary posterior multimodality,
skewness, or hard global interaction structure.  It defines one deterministic
Gaussian-surrogate lane and makes its approximation boundary explicit.

\paragraph{How to read this note.}
Sections~\ref{sec:problem-and-lane}--\ref{sec:algorithm-one-pass} give the
problem statement, the exact-versus-approximate object distinction, and the full
runtime path.  Sections~\ref{sec:sparse-grid-construction} and
\ref{sec:value-path} define the fixed cloud and the value recursion.
Sections~\ref{sec:same-scalar-contract} and \ref{sec:gradient} define the saved
branch and the analytical gradient.  Section~\ref{sec:fd-check} defines the
finite-difference audit, and Section~\ref{sec:validation} states what evidence
supports implementation correctness versus what evidence supports only
approximate scientific usefulness.

\section{Problem, Exact Target, and Selected Approximation Lane}
\label{sec:problem-and-lane}

The scientific problem is high-dimensional nonlinear filtering.  In a general
state-space model, the exact Bayesian filter propagates a full conditional law
through time.  Once the transition or observation map is nonlinear, that law is
usually no longer representable by a finite list of moments, and in high
dimension exact function-valued propagation is rarely a practical computational
object.

The question addressed here is therefore narrower than exact filtering: what
deterministic approximation lane is transparent about its compromises,
mathematically coherent, and still implementable with a reproducible value and
analytical gradient?

The lane selected here is \emph{fixed sparse-grid quadrature filtering}
(FixedSGQF).  The method makes three deliberate approximation decisions:
\begin{enumerate}[leftmargin=0.3in]
  \item it carries one Gaussian surrogate \(\N(m_t,P_t)\) instead of the full
  filtering density,
  \item it evaluates the nonlinear moments needed by that Gaussian projection
  with a deterministic sparse-grid quadrature rule in standardized Gaussian
  coordinates, and
  \item it freezes the cloud, merge rule, covariance-factor branch, and veto
  policy so that the reported gradient differentiates the same deterministic
  scalar that the reported value routine evaluates.
\end{enumerate}

This note chooses FixedSGQF because it defines an auditable deterministic
Gaussian-surrogate scalar, an analytical derivative of that same scalar on a
fixed branch, and a moment engine less explosive than dense tensor-product
Gaussian quadrature in the intended high-dimensional regime
\cite{jia2012sparsegrid,singh2018adaptive}.  It is not a substitute for richer
non-Gaussian density-approximation lanes such as the Zhao--Cui squared
tensor-train construction \cite{zhao2024ttsequential}; it is the deterministic
Gaussian-surrogate lane in a broader program.

\section{What the Algorithm Computes}
\label{sec:what-computes}

This section is the scope statement an implementer should keep in mind while
reading every later formula.

\paragraph{Exact target.}
The scientific target is the exact nonlinear filtering law
\begin{equation}
  p_\theta(x_t\mid y_{1:t}).
  \label{eq:exact-filtering-law}
\end{equation}
That is a full probability law on \(\R^{n_x}\).

\paragraph{Approximate carried object.}
FixedSGQF does \emph{not} store that full law.  It stores only one Gaussian
surrogate,
\begin{equation}
  X_t\mid y_{1:t}
  \approx
  \N(m_t,P_t).
  \label{eq:carried-gaussian}
\end{equation}

\paragraph{Accumulated deterministic scalar.}
At each time step the method forms a Gaussian innovation log-likelihood
increment,
\begin{equation}
  \ellhat_t^{\rm FSGQ}
  =
  -\frac12\left\{\log\det S_t + v_t^\top S_t^{-1}v_t + n_y\log(2\pi)\right\},
  \label{eq:increment-loglik}
\end{equation}
and then accumulates
\begin{equation}
  \ellhat_T^{\rm FSGQ}(\theta)=\sum_{t=1}^T \ellhat_t^{\rm FSGQ}(\theta).
  \label{eq:total-loglik}
\end{equation}
The analytical gradient in this note is the derivative of the same scalar in
\eqref{eq:total-loglik}, not the derivative of an adaptive grid-selection
procedure.

\paragraph{What is discarded.}
The approximation in \eqref{eq:carried-gaussian} discards general non-Gaussian
posterior shape.  Sparse-grid exactness for selected Gaussian moments is not the
same as exact posterior representation.

\paragraph{A scalar warning cell.}
Let
\begin{equation}
  X\sim\N(0,P),
  \qquad
  Y=X^2+\varepsilon,
  \qquad
  \varepsilon\sim\N(0,R).
  \label{eq:quadratic-cell}
\end{equation}
For \(y>0\) and small \(R\), the exact posterior density can have two separated
lobes near \(\pm\sqrt y\).  A Gaussian-projection filter can still compute useful
moments of \(X^2\), but after the update it stores only one Gaussian.  That is
not a bug in quadrature.  It is the design choice in
\eqref{eq:carried-gaussian}.

\section{Object Map and Core Terminology}
\label{sec:object-map}

The table below is the object inventory an implementer should keep open while
coding.

\begin{center}
\small
\begin{tabular}{p{0.22\linewidth}p{0.45\linewidth}p{0.21\linewidth}}
\toprule
symbol & meaning & shape \\
\midrule
\(\xi^{(r)}\) & standardized sparse-grid node & \(n_x\) \\
\(w_r\) & accumulated signed weight after merge & scalar \\
\(m_t,P_t\) & carried Gaussian after update & \(n_x\), \(n_x\times n_x\) \\
\(m_t^-,P_t^-\) & predictive Gaussian before update & \(n_x\), \(n_x\times n_x\) \\
\(C_t,C_t^-\) & Cholesky factors of \(P_t\) and \(P_t^-\) & \(n_x\times n_x\) \\
\(\chi_{t-1}^{(r)}\) & placed previous-state cloud point & \(n_x\) \\
\(a_t^{(r)}\) & transition image of \(\chi_{t-1}^{(r)}\) & \(n_x\) \\
\(\chi_t^{(r)}\) & placed predictive cloud point & \(n_x\) \\
\(z_t^{(r)}\) & observation image of \(\chi_t^{(r)}\) & \(n_y\) \\
\(\bar z_t\) & predicted observation mean & \(n_y\) \\
\(S_t\) & innovation covariance & \(n_y\times n_y\) \\
\(C_{xz,t}\) & state-observation cross covariance & \(n_x\times n_y\) \\
\(v_t\) & innovation residual \(y_t-\bar z_t\) & \(n_y\) \\
\(K_t\) & Gaussian-projection gain & \(n_x\times n_y\) \\
\(\ellhat_t^{\rm FSGQ}\) & one-step deterministic approximate log likelihood & scalar \\
\(\dot{\square}\) & derivative with respect to one parameter component & matching object \\
\bottomrule
\end{tabular}
\end{center}

\paragraph{Exact filtering law.}
The full conditional law \(p_\theta(x_t\mid y_{1:t})\).

\paragraph{Carried Gaussian surrogate.}
The pair \((m_t,P_t)\) used in place of the full law.

\paragraph{Standardized sparse-grid cloud.}
A fixed list \(\{(\xi^{(r)},w_r)\}_{r=1}^M\) in standard normal coordinates.
The cloud is fixed before value evaluation.

\paragraph{Physical cloud points.}
Points obtained by mapping standardized nodes through the current mean and
covariance factor, for example \(x=m+C\xi\).

\paragraph{Fixed branch.}
The saved structural choices under which the value and derivative are both
evaluated: cloud, merge rule, ordering, factor convention, veto thresholds, and
other declared policies.

\paragraph{Same-scalar derivative.}
The derivative of the deterministic scalar actually evaluated on the saved
branch, not the derivative of a changing adaptive algorithm.

\paragraph{Branch veto.}
A declared failure exit, such as a predictive covariance or innovation
covariance failing the positive-definite branch check.

\paragraph{Implementation conventions.}
This note fixes three conventions report-wide:
\begin{enumerate}[leftmargin=0.3in]
  \item every covariance factor is the lower-triangular Cholesky factor with
  positive diagonal, written \(C(P)=\chol(P)\), whenever the declared branch is
  valid;
  \item all expressions written with \(S_t^{-1}\) are implemented by linear
  solves against the Cholesky factor of \(S_t\), not by forming an explicit
  inverse;
  \item if floating-point accumulation yields a covariance that is not exactly
  symmetric, the implementation first symmetrizes by
  \((P+P^\top)/2\) and then either accepts the declared branch or returns a veto;
  it does not silently repair the scalar with hidden jitter, clipping, or floor
  operations.
\end{enumerate}

\section{FixedSGQF in One Pass}
\label{sec:algorithm-one-pass}

A junior implementer should be able to say what the method does after one read:
\begin{enumerate}[leftmargin=0.3in]
  \item Build a fixed sparse-grid cloud in standardized Gaussian coordinates.
  \item Start from an initial Gaussian surrogate \((m_0,P_0)\).
  \item At each time step, place the previous cloud using the current factor
  \(C_{t-1}\), pass those points through the nonlinear transition map, and form
  predictive moments \((m_t^-,P_t^-)\).
  \item If the predictive covariance fails the declared branch, stop and report
  a predictive-branch failure.
  \item Place the predictive cloud using \(C_t^-\), pass those points through the
  nonlinear observation map, and form \(\bar z_t\), \(S_t\), and
  \(C_{xz,t}\).
  \item If the innovation covariance fails the declared branch, stop and report
  an innovation-branch failure.
  \item Form the one-step Gaussian innovation log likelihood
  \(\ellhat_t^{\rm FSGQ}\), update the running total
  \(\ellhat_T^{\rm FSGQ}\), and update the carried Gaussian
  \((m_t,P_t)\).
  \item If the updated carried covariance fails the declared branch, stop and
  report a carried-covariance failure.
  \item If gradients are requested, propagate the same fixed-branch derivative
  recursion using the same cloud and the same branch policies.
\end{enumerate}

The approximation ladder is therefore:
\begin{equation}
\begin{aligned}
  \text{exact target} &\Longrightarrow p_\theta(x_t\mid y_{1:t}), \\
  \text{carried surrogate} &\Longrightarrow \N(m_t,P_t), \\
  \text{moment engine} &\Longrightarrow \text{fixed sparse-grid quadrature}, \\
  \text{reported scalar} &\Longrightarrow \ellhat_T^{\rm FSGQ}(\theta), \\
  \text{reported derivative} &\Longrightarrow \nabla_\theta \ellhat_T^{\rm FSGQ}(\theta;\branch).
\end{aligned}
\label{eq:approximation-ladder}
\end{equation}

\section{State-Space Model and Gaussian Projection}
\label{sec:ssm-gaussian-projection}

Use the additive-noise state-space model
\begin{align}
  X_t &= f_\theta(X_{t-1}) + \eta_t,
  &
  \eta_t &\sim \N(0,Q_\theta),
  \label{eq:state-model}\\
  Y_t &= h_\theta(X_t) + \varepsilon_t,
  &
  \varepsilon_t &\sim \N(0,R_\theta),
  \label{eq:observation-model}
\end{align}
with \(X_t\in\R^{n_x}\), \(Y_t\in\R^{n_y}\), and parameter
\(\theta\in\R^p\).

The exact Bayesian prediction and update are
\begin{align}
  p_\theta(x_t\mid y_{1:t-1})
  &=
  \int p_\theta(x_t\mid x_{t-1}) p_\theta(x_{t-1}\mid y_{1:t-1})\dd x_{t-1},
  \label{eq:exact-prediction}\\
  p_\theta(x_t\mid y_{1:t})
  &=
  \frac{g_\theta(y_t\mid x_t)p_\theta(x_t\mid y_{1:t-1})}
  {\int g_\theta(y_t\mid x_t)p_\theta(x_t\mid y_{1:t-1})\dd x_t}.
  \label{eq:exact-update}
\end{align}
FixedSGQF does not close these exact densities; it performs a Gaussian moment
projection.

Assume the predictive state is represented by
\begin{equation}
  X_t^-\sim\N(m_t^-,P_t^-).
  \label{eq:predictive-gaussian}
\end{equation}
Let \(Z_t=h_\theta(X_t^-)\).  The Gaussian projection uses
\begin{align}
  \bar z_t &= \E[Z_t],
  \label{eq:zbar}\\
  S_t &= R_\theta + \E[(Z_t-\bar z_t)(Z_t-\bar z_t)^\top],
  \label{eq:S-def}\\
  C_{xz,t} &= \E[(X_t^- - m_t^-)(Z_t-\bar z_t)^\top].
  \label{eq:Cxz-def}
\end{align}
Given these moments, define
\begin{align}
  v_t &= y_t-\bar z_t,
  &
  K_t &= C_{xz,t}S_t^{-1},
  \label{eq:innovation-and-gain}\\
  m_t &= m_t^- + K_t v_t,
  &
  P_t &= P_t^- - K_tS_tK_t^\top.
  \label{eq:gaussian-projection-update}
\end{align}
The sparse-grid machinery answers how to approximate the expectations in
\eqref{eq:zbar}--\eqref{eq:Cxz-def}; it does not change the fact that the
carried object is \((m_t,P_t)\).

\section{Sparse-Grid Construction in Standardized Coordinates}
\label{sec:sparse-grid-construction}

Let \(\xi\sim\N(0,I_{n_x})\).  The sparse-grid cloud is built once in
standardized coordinates and then mapped into physical coordinates by the
current Gaussian factor.

For a one-dimensional standard-normal quadrature rule at level \(\ell\), write
\begin{equation}
  I_\ell[\psi]
  =
  \sum_{a=1}^{s_\ell}\omega_{\ell,a}\psi(\zeta_{\ell,a})
  \approx
  \int_\R \psi(\xi)\phi(\xi)\dd\xi.
  \label{eq:univariate-rule}
\end{equation}
For a multi-index \(\mathbf i=(i_1,\ldots,i_b)\), the tensor rule is
\begin{equation}
  I_{\mathbf i}[F] = (I_{i_1}\otimes\cdots\otimes I_{i_b})[F].
  \label{eq:tensor-rule}
\end{equation}

Following \citet[Eq.~26--29]{jia2012sparsegrid}, define the active band
\begin{equation}
  \mathcal A_{b,L}
  =
  \{\mathbf i\in\mathbb N^b: L\le |\mathbf i|\le L+b-1\},
  \qquad
  |\mathbf i|=\sum_{j=1}^b i_j,
  \label{eq:active-band}
\end{equation}
and the Smolyak coefficient
\begin{equation}
  c_{\mathbf i}
  =
  (-1)^{L+b-1-|\mathbf i|}
  \binom{b-1}{L+b-1-|\mathbf i|}.
  \label{eq:smolyak-coeff}
\end{equation}
Then
\begin{equation}
  A_{b,L}[F]
  =
  \sum_{\mathbf i\in\mathcal A_{b,L}} c_{\mathbf i}I_{\mathbf i}[F].
  \label{eq:smolyak-rule}
\end{equation}

Each tensor point contributes a standardized node and signed weight.  Duplicate
standardized nodes must be merged by accumulating their signed contributions.
The final fixed cloud is
\begin{equation}
  \mathcal C_{b,L} = \{(\xi^{(r)},w_r)\}_{r=1}^{M_{b,L}}.
  \label{eq:fixed-cloud}
\end{equation}

\paragraph{Cloud-construction contract.}
The fixed cloud is part of the saved scalar.  The univariate rule family, merge
tolerance, zero-weight threshold, enumeration order, and final sort order all
belong to the target.  Two implementations that use the same level but different
merge or ordering policies have not necessarily evaluated the same scalar.

\subsection{Pseudocode: build the fixed sparse-grid cloud}
\label{subsec:cloud-pseudocode}

\fbox{\begin{minipage}{0.96\linewidth}
\textbf{Routine: \texttt{build\_fixed\_sparse\_grid\_cloud}}\\[0.2em]
\textbf{Inputs:} dimension \(b\), level \(L\), one-dimensional rule family
\(\{I_\ell\}\), merge tolerance \(\tau_{\rm merge}\), zero-weight threshold
\(\tau_{\rm zero}\).\\
\textbf{Outputs:} merged standardized cloud
\(\mathcal C_{b,L}=\{(\xi^{(r)},w_r)\}_{r=1}^{M_{b,L}}\), plus diagnostics.
\begin{enumerate}[leftmargin=0.28in]
  \item Enumerate all multi-indices \(\mathbf i\in\mathcal A_{b,L}\) in
  lexicographic order.
  \item For each \(\mathbf i\), enumerate local tensor indices \(\alpha\) in
  lexicographic order.
  \item Form the candidate standardized node
  \(\xi^{\mathbf i,\alpha}\) and signed contribution
  \(c_{\mathbf i}w^{\mathbf i,\alpha}\).
  \item Search the current dictionary for an existing node \(\xi\) satisfying
  \(\|\xi^{\mathbf i,\alpha}-\xi\|_\infty\le \tau_{\rm merge}\).
  \item If found, add the signed contribution to the stored weight for that
  node; otherwise create a new dictionary entry.
  \item After all contributions have been processed, remove entries with
  accumulated absolute weight below \(\tau_{\rm zero}\).
  \item Sort the surviving nodes lexicographically and store them in that
  order.
  \item Report at least the weight-total diagnostic
  \(\sum_r w_r\), the number of surviving nodes, and the number of negative
  accumulated weights.
\end{enumerate}
\end{minipage}}

\subsection{A toy cloud oracle}
\label{subsec:toy-cloud}

Take \(b=2\) and \(L=2\).  Use the level-1 and level-2 one-dimensional rules
\begin{equation}
  I_1:(0,1),
  \qquad
  I_2:\left\{(-\sqrt3,1/6),(0,2/3),(\sqrt3,1/6)\right\}.
  \label{eq:toy-rules}
\end{equation}
The active band contains
\begin{equation}
  (1,1),\qquad (1,2),\qquad (2,1),
  \label{eq:toy-indices}
\end{equation}
with coefficients
\begin{equation}
  c_{(1,1)}=-1,
  \qquad
  c_{(1,2)}=1,
  \qquad
  c_{(2,1)}=1.
  \label{eq:toy-coeffs}
\end{equation}
After duplicate merge, the cloud is
\begin{center}
\begin{tabular}{ccc}
\toprule
node \(\xi\) & accumulated weight & source components \\
\midrule
\((0,0)\) & \(1/3\) & \((1,1),(1,2),(2,1)\) \\
\(( -\sqrt3,0 )\) & \(1/6\) & \((2,1)\) \\
\(( \sqrt3,0 )\) & \(1/6\) & \((2,1)\) \\
\((0,-\sqrt3)\) & \(1/6\) & \((1,2)\) \\
\((0,\sqrt3)\) & \(1/6\) & \((1,2)\) \\
\bottomrule
\end{tabular}
\end{center}
The weights sum to one, and any implementation should reproduce this merged
cloud exactly up to the declared merge and zero-weight tolerances.

\section{Value Path: Filtering and Deterministic Likelihood}
\label{sec:value-path}

Fix a merged standardized cloud
\begin{equation}
  \mathcal C=\{(\xi^{(r)},w_r)\}_{r=1}^M,
  \qquad
  \xi^{(r)}\in\R^{n_x},
  \quad
  w_r\in\R.
  \label{eq:cloud-for-filter}
\end{equation}
The cloud is fixed before likelihood evaluation.

\subsection{Prediction stage}

Let \(P_{t-1}=C_{t-1}C_{t-1}^\top\) on the declared Cholesky branch.  Place the
previous-state cloud,
\begin{equation}
  \chi_{t-1}^{(r)} = m_{t-1}+C_{t-1}\xi^{(r)},
  \label{eq:pred-placed-points}
\end{equation}
push those points through the transition,
\begin{equation}
  a_t^{(r)} = f_\theta(\chi_{t-1}^{(r)}),
  \label{eq:transition-points}
\end{equation}
and compute the predictive moments
\begin{align}
  m_t^- &= \sum_{r=1}^M w_r a_t^{(r)},
  \label{eq:pred-mean}\\
  P_t^- &= Q_\theta + \sum_{r=1}^M w_r
  (a_t^{(r)}-m_t^-)(a_t^{(r)}-m_t^-)^\top.
  \label{eq:pred-cov}
\end{align}
After \eqref{eq:pred-cov}, the implementation symmetrizes by
\((P_t^-+P_t^{-\top})/2\) and either accepts the Cholesky branch or returns a
predictive-branch veto.

\subsection{Observation stage}

Factor \(P_t^-=C_t^-(C_t^-)^\top\) on the declared branch and place the
predictive cloud,
\begin{equation}
  \chi_t^{(r)} = m_t^- + C_t^-\xi^{(r)}.
  \label{eq:obs-placed-points}
\end{equation}
Evaluate the observation map,
\begin{equation}
  z_t^{(r)} = h_\theta(\chi_t^{(r)}),
  \label{eq:obs-values}
\end{equation}
and compute
\begin{align}
  \bar z_t &= \sum_{r=1}^M w_r z_t^{(r)},
  \label{eq:obs-mean}\\
  \Delta_t^{(r)} &= z_t^{(r)}-\bar z_t,
  \label{eq:obs-centered}\\
  S_t &= R_\theta + \sum_{r=1}^M w_r \Delta_t^{(r)}\Delta_t^{(r)\top},
  \label{eq:obs-cov}\\
  C_{xz,t} &= \sum_{r=1}^M w_r
  (\chi_t^{(r)}-m_t^-)\Delta_t^{(r)\top}.
  \label{eq:cross-cov}
\end{align}
If the symmetrized innovation covariance fails the declared branch, the value
path returns an innovation-branch veto.

If it passes, define
\begin{align}
  v_t &= y_t-\bar z_t,
  \label{eq:innovation}\\
  \ellhat_t^{\rm FSGQ}
  &=
  -\frac12\left\{\log\det S_t + v_t^\top S_t^{-1}v_t + n_y\log(2\pi)\right\},
  \label{eq:value-loglik}\\
  K_t &= C_{xz,t}S_t^{-1},
  \label{eq:kalman-gain}\\
  m_t &= m_t^- + K_t v_t,
  \label{eq:filter-mean}\\
  P_t &= P_t^- - K_tS_tK_t^\top.
  \label{eq:filter-cov}
\end{align}
After the update, the implementation again symmetrizes \(P_t\) and either
accepts the carried branch or returns a carried-covariance veto.

\subsection{Pseudocode: one-step value recursion}
\label{subsec:value-pseudocode}

\fbox{\begin{minipage}{0.96\linewidth}
\textbf{Routine: \texttt{fixed\_sgqf\_filter\_value\_step}}\\[0.2em]
\textbf{Inputs:} previous mean/covariance \((m_{t-1},P_{t-1})\), current
observation \(y_t\), model objects \(f_\theta,h_\theta,Q_\theta,R_\theta\),
fixed cloud \(\mathcal C\), branch thresholds, factor convention.\\
\textbf{Outputs:} one-step increment \(\ellhat_t^{\rm FSGQ}\), updated pair
\((m_t,P_t)\), and either acceptance diagnostics or a failure record.
\begin{enumerate}[leftmargin=0.28in]
  \item Symmetrize \(P_{t-1}\) and factor it as
  \(P_{t-1}=C_{t-1}C_{t-1}^\top\).  If factorization fails, return a
  factor-branch failure.
  \item Place previous-state points \(\chi_{t-1}^{(r)}\) and compute transition
  images \(a_t^{(r)}\).
  \item Compute \(m_t^-\) and \(P_t^-\), then symmetrize \(P_t^-\).  If the
  predictive branch fails, return a predictive-covariance failure.
  \item Factor \(P_t^-\) as \(C_t^-C_t^{-\top}\) and place predictive points
  \(\chi_t^{(r)}\).
  \item Compute observation images \(z_t^{(r)}\), then
  \(\bar z_t\), \(S_t\), and \(C_{xz,t}\).  Symmetrize \(S_t\).  If the
  innovation branch fails, return an innovation-covariance failure.
  \item Form \(v_t\), solve against \(S_t\), compute
  \(\ellhat_t^{\rm FSGQ}\), and update the running scalar.
  \item Compute \(K_t\), \(m_t\), and \(P_t\).  Symmetrize \(P_t\).  If the
  carried branch fails, return a carried-covariance failure.
  \item Return \(\ellhat_t^{\rm FSGQ}\), \((m_t,P_t)\), and diagnostics.
\end{enumerate}
\end{minipage}}

\section{Worked One-Step Numerical Oracle}
\label{sec:worked-example}

This section is a compact end-to-end oracle for both the value path and the
gradient path.

Use the one-dimensional model
\begin{align}
  X_0 &\sim \N(m_0,P_0),
  &m_0&=0.5,
  &P_0&=1,
  \label{eq:example-initial}\\
  X_1 &= f_\beta(X_0)+\eta_1,
  &f_\beta(x)&=\beta x,
  &\eta_1&\sim\N(0,Q),
  \label{eq:example-transition}\\
  Y_1 &= h(X_1)+\varepsilon_1,
  &h(x)&=x^2,
  &\varepsilon_1&\sim\N(0,R),
  \label{eq:example-observation}
\end{align}
with
\begin{equation}
  Q=0.25,
  \qquad
  R=1,
  \qquad
  y_1=2,
  \qquad
  \beta=1.
  \label{eq:example-QRy}
\end{equation}
Use the three-point standard-normal cloud
\begin{equation}
  \mathcal C_{1,2}
  =
  \left\{\left(-\sqrt3,\frac16\right),\left(0,\frac23\right),
  \left(\sqrt3,\frac16\right)\right\}.
  \label{eq:example-cloud}
\end{equation}
Then the one-step value objects are
\begin{align}
  m_1^- &= 0.5,
  &
  P_1^- &= 1.25 = \frac54,
  \label{eq:example-pred-moments}\\
  \bar z_1 &= 1.5 = \frac32,
  &
  S_1 &= 5.375 = \frac{43}{8},
  \label{eq:example-observation-moments}\\
  C_{xz,1} &= 1.25 = \frac54,
  &
  v_1 &= 0.5 = \frac12,
  \label{eq:example-cross-and-innovation}\\
  K_1 &= \frac{10}{43}\approx 0.2325581395,
  &
  m_1 &= \frac{53}{86}\approx 0.6162790698,
  \label{eq:example-gain-and-mean}\\
  P_1 &= \frac{165}{172}\approx 0.9593023256,
  &
  \ellhat_1^{\rm FSGQ} &\approx -1.7830736342.
  \label{eq:example-final-values}
\end{align}
An implementation should reproduce these values up to declared floating-point
tolerance before attempting larger cases.

\section{Saved Scalar and Same-Scalar Contract}
\label{sec:same-scalar-contract}

The value and derivative paths are meaningful only if they refer to the same
saved scalar.

\paragraph{Plain-language rule.}
Value, gradient, and finite differences must use the same fixed cloud, the same
merge and ordering rules, the same factor family, the same symmetrize-then-veto
policy, and the same preprocessing and initial-condition policies.  The
parameter-dependent numerical objects are recomputed on that same branch; the
structural route is what must remain fixed.

Formally, a saved branch is the tuple
\begin{equation}
\begin{aligned}
  \branch
  =
  (&y_{1:T},\ \text{observation preprocessing},\ m_0(\theta),P_0(\theta),\dot m_0,\dot P_0,\\
  &\mathcal C,\ \text{one-dimensional rules},\ \text{merge tolerance},\ \text{zero-weight rule},\ \text{ordering},\\
  &C(P)=\chol(P),\ \text{symmetrize-then-veto rule},\ \epsilon_S,\epsilon_P,\\
  &f_\theta,h_\theta,Q_\theta,R_\theta,
  D_xf_\theta,\partial_i f_\theta,D_xh_\theta,\partial_i h_\theta,
  \dot Q_\theta,\dot R_\theta).
\end{aligned}
\label{eq:branch-tuple}
\end{equation}
The fixed scalar is
\begin{equation}
  \ellhat_T^{\rm FSGQ}(\theta;\branch)
  =
  \sum_{t=1}^T \ellhat_t^{\rm FSGQ}(\theta;\branch).
  \label{eq:fixed-scalar-branch}
\end{equation}

\paragraph{Frozen versus variable objects.}
The frozen structural objects are the cloud, merge rule, zero-weight rule,
ordering, factor family, thresholds, preprocessing policy, and
initial-condition policy.  The parameter-dependent numerical objects that are
recomputed on that fixed branch include
\(m_t,P_t,C_t,m_t^-,P_t^-,C_t^-,\bar z_t,S_t,C_{xz,t},K_t\), and
\(\ellhat_t^{\rm FSGQ}\).

\paragraph{Operational branch identity for finite differences.}
A finite-difference row is branch-valid only if both perturbed evaluations reuse
the same structural metadata and follow the same realized acceptance path through
each likelihood contribution.  If one perturbation triggers a different first
failure location, a different accepted-prefix pattern, or a different structural
policy, the row is branch-invalid and is not interpreted as derivative evidence.

\section{Analytical Gradient of the Fixed Scalar}
\label{sec:gradient}

\paragraph{Why this derivative is nontrivial.}
The one-step log likelihood at time \(t\) depends on objects built by earlier
steps of the filter.  The derivative must therefore propagate through the
previous carried Gaussian, the covariance factors used to place points, the
nonlinear transition and observation maps, the innovation covariance, and the
posterior update.  The branch must be fixed first because otherwise the value and
derivative would refer to different targets.

\paragraph{Score objects versus propagation objects.}
For the current one-step score, the essential derivative objects are
\(\dot v_t\) and \(\dot S_t\).  The derivative of the cross covariance
\(\dot C_{xz,t}\) is not needed to close the current score, but it is needed to
form \(\dot K_t\), and hence to propagate \(\dot m_t\) and \(\dot P_t\) to the
next time step.

Differentiate with respect to one parameter component \(\theta_i\).  A dot
denotes \(\partial_i\).

The derivative dependency chain is
\begin{equation}
\begin{aligned}
  \theta
  &\longmapsto (m_{t-1},P_{t-1})
  \longmapsto C_{t-1}
  \longmapsto \chi_{t-1}^{(r)}
  \longmapsto a_t^{(r)}
  \longmapsto (m_t^-,P_t^-) \\
  &\longmapsto C_t^-
  \longmapsto \chi_t^{(r)}
  \longmapsto z_t^{(r)}
  \longmapsto (\bar z_t,S_t,C_{xz,t})
  \longmapsto (v_t,K_t,\ellhat_t,m_t,P_t).
\end{aligned}
\label{eq:gradient-chain}
\end{equation}

\subsection{Cholesky-branch derivative}
\label{subsec:chol-derivative}

If \(P=LL^\top\) with positive diagonal and \(C(P)=L\), define
\begin{equation}
  A = L^{-1}\dot P L^{-\top}.
  \label{eq:chol-A}
\end{equation}
Let \(\Phi(A)\) be the lower-triangular operator
\begin{equation}
  \Phi(A)_{jk}
  =
  \begin{cases}
    A_{jk}, & j>k,\\
    A_{jj}/2, & j=k,\\
    0, & j<k.
  \end{cases}
  \label{eq:chol-Phi}
\end{equation}
Then
\begin{equation}
  \dot L = L\Phi(A).
  \label{eq:chol-derivative}
\end{equation}
This derivative matters operationally because every placed cloud point depends on
the covariance factor.

\subsection{Prediction sensitivities}

Differentiate the placed previous-state points,
\begin{equation}
  \dot\chi_{t-1}^{(r)} = \dot m_{t-1}+\dot C_{t-1}\xi^{(r)},
  \label{eq:dot-pred-place}
\end{equation}
and the transition images,
\begin{equation}
  \dot a_t^{(r)} = D_xf_\theta(\chi_{t-1}^{(r)})\dot\chi_{t-1}^{(r)} + \partial_i f_\theta(\chi_{t-1}^{(r)}).
  \label{eq:dot-transition}
\end{equation}
Then
\begin{align}
  \dot m_t^- &= \sum_{r=1}^M w_r\dot a_t^{(r)},
  \label{eq:dot-pred-mean}\\
  d_t^{(r)} &= a_t^{(r)}-m_t^-,
  \qquad
  \dot d_t^{(r)} = \dot a_t^{(r)}-\dot m_t^-,
  \label{eq:pred-centered}\\
  \dot P_t^- &= \dot Q_\theta + \sum_{r=1}^M w_r\left\{
  \dot d_t^{(r)}d_t^{(r)\top} + d_t^{(r)}\dot d_t^{(r)\top}
  \right\}.
  \label{eq:dot-pred-cov}
\end{align}

\subsection{Observation sensitivities}

Differentiate the predictive placed points,
\begin{equation}
  \dot\chi_t^{(r)} = \dot m_t^- + \dot C_t^-\xi^{(r)},
  \label{eq:dot-obs-place}
\end{equation}
then the observation images,
\begin{equation}
  \dot z_t^{(r)} = D_xh_\theta(\chi_t^{(r)})\dot\chi_t^{(r)} + \partial_i h_\theta(\chi_t^{(r)}),
  \label{eq:dot-obs-value}
\end{equation}
and then
\begin{align}
  \dot{\bar z}_t &= \sum_{r=1}^M w_r\dot z_t^{(r)},
  \label{eq:dot-zbar}\\
  \dot v_t &= -\dot{\bar z}_t,
  \label{eq:dot-v}\\
  \dot\Delta_t^{(r)} &= \dot z_t^{(r)}-\dot{\bar z}_t,
  \label{eq:dot-Delta}\\
  \dot S_t &= \dot R_\theta + \sum_{r=1}^M w_r\left\{
  \dot\Delta_t^{(r)}\Delta_t^{(r)\top} + \Delta_t^{(r)}\dot\Delta_t^{(r)\top}
  \right\},
  \label{eq:dot-S}\\
  \dot C_{xz,t} &= \sum_{r=1}^M w_r\left\{
  (\dot\chi_t^{(r)}-\dot m_t^-)\Delta_t^{(r)\top} +
  (\chi_t^{(r)}-m_t^-)\dot\Delta_t^{(r)\top}
  \right\}.
  \label{eq:dot-Cxz}
\end{align}

\subsection{Innovation score}

\begin{proposition}[Fixed Gaussian innovation score]
\label{prop:innovation-score}
Assume \(S_t\) is positive definite and the saved branch is differentiable.
Let \(u_t\) solve
\begin{equation}
  S_tu_t=v_t.
  \label{eq:score-solve}
\end{equation}
Then the derivative of \eqref{eq:value-loglik} is
\begin{equation}
  \dot\ellhat_t^{\rm FSGQ}
  =
  -\frac12\left[\tr(S_t^{-1}\dot S_t) + 2\dot v_t^\top u_t - u_t^\top \dot S_tu_t\right].
  \label{eq:score-formula}
\end{equation}
\end{proposition}

\begin{proof}
Use \(\dd\log\det S = \tr(S^{-1}\dd S)\) and
\(\dd S^{-1} = -S^{-1}(\dd S)S^{-1}\).  With \(u=S^{-1}v\),
\begin{equation}
  \dd(v^\top S^{-1}v) = 2(\dd v)^\top u - u^\top(\dd S)u.
  \label{eq:d-quadratic}
\end{equation}
Substitute \(\dd v=\dot v_t\) and \(\dd S=\dot S_t\) into
\eqref{eq:value-loglik}.
\end{proof}

\subsection{Posterior sensitivity propagation}

Differentiate the gain,
\begin{equation}
  \dot K_t = \dot C_{xz,t}S_t^{-1} - K_t\dot S_tS_t^{-1},
  \label{eq:dot-K}
\end{equation}
then propagate
\begin{align}
  \dot m_t &= \dot m_t^- + \dot K_t v_t + K_t\dot v_t,
  \label{eq:dot-m}\\
  \dot P_t &= \dot P_t^- - \dot K_tS_tK_t^\top - K_t\dot S_tK_t^\top - K_tS_t\dot K_t^\top.
  \label{eq:dot-P}
\end{align}
The structural summary is: \(\dot v_t\) and \(\dot S_t\) close the current
score, while \(\dot C_{xz,t}\), \(\dot K_t\), \(\dot m_t\), and
\(\dot P_t\) exist because the derivative state must be carried forward in
time.

\subsection{Pseudocode: one-step gradient recursion}
\label{subsec:gradient-pseudocode}

\fbox{\begin{minipage}{0.96\linewidth}
\textbf{Routine: \texttt{fixed\_sgqf\_filter\_gradient\_step}}\\[0.2em]
\textbf{Inputs:} all value-step inputs, plus derivative objects
\((\dot m_{t-1},\dot P_{t-1})\), first-derivative model routines, and the same
saved branch.\\
\textbf{Outputs:} one-step derivative contribution
\(\dot\ellhat_t^{\rm FSGQ}\), updated derivative state
\((\dot m_t,\dot P_t)\), and either acceptance diagnostics or the same failure
record as the value path.
\begin{enumerate}[leftmargin=0.28in]
  \item Reuse the same accepted value-path branch through previous factor,
  predictive factor, and innovation covariance.
  \item Compute \(\dot C_{t-1}\), then \(\dot\chi_{t-1}^{(r)}\) and
  \(\dot a_t^{(r)}\).
  \item Compute \(\dot m_t^-\) and \(\dot P_t^-\), then \(\dot C_t^-\).
  \item Compute \(\dot\chi_t^{(r)}\), \(\dot z_t^{(r)}\),
  \(\dot{\bar z}_t\), \(\dot v_t\), \(\dot S_t\), and
  \(\dot C_{xz,t}\).
  \item Solve \(S_tu_t=v_t\) on the accepted innovation branch and evaluate the
  one-step score derivative with \eqref{eq:score-formula}.
  \item Compute \(\dot K_t\), then propagate \(\dot m_t\) and \(\dot P_t\).
  \item Return the derivative contribution, the updated derivative state, and
  diagnostics.
\end{enumerate}
\end{minipage}}

\section{Finite-Difference Same-Scalar Check}
\label{sec:fd-check}

Finite differences test whether the implementation differentiates the same saved
scalar that the value path evaluates.

For parameter component \(i\), define the central difference
\begin{equation}
  D_i(h)
  =
  \frac{\ellhat_T^{\rm FSGQ}(\theta+he_i;\branch)-
        \ellhat_T^{\rm FSGQ}(\theta-he_i;\branch)}{2h}.
  \label{eq:central-diff}
\end{equation}
Use a step ladder such as
\begin{equation}
  h\in\{10^{-1},10^{-2},10^{-3},10^{-4},10^{-5}\}\max\{1,|\theta_i|\}.
  \label{eq:fd-ladder}
\end{equation}
A row is branch-valid only if both perturbed evaluations reuse the same saved
structural metadata and match the base evaluation in their realized
accepted/failure stage-time pattern.  If either side changes branch or returns a
different veto pattern, the row is branch-invalid and is not interpreted as
score evidence.

For valid rows, compare the analytical score \(G_i\) against \(D_i(h)\) by
\begin{equation}
  e_i(h) = \frac{|D_i(h)-G_i|}{\max\{1,|D_i(h)|,|G_i|\}}.
  \label{eq:fd-relative-error}
\end{equation}
A practical diagnostic passes when at least two consecutive valid rows satisfy
\(e_i(h)\le 10^{-5}\) and the smaller-step row is no worse than ten times the
larger-step row.

\subsection{Pseudocode: same-scalar finite-difference audit}
\label{subsec:fd-pseudocode}

\fbox{\begin{minipage}{0.96\linewidth}
\textbf{Routine: \texttt{same\_scalar\_fd\_check}}\\[0.2em]
\textbf{Inputs:} parameter component \(i\), step ladder, base parameter value,
saved branch \(\branch\), value routine, analytical score \(G_i\).\\
\textbf{Outputs:} finite-difference table with branch-valid flags and relative
errors.
\begin{enumerate}[leftmargin=0.28in]
  \item For each step size \(h\), evaluate the value routine at
  \(\theta+he_i\) and \(\theta-he_i\) using the same saved structural branch.
  \item Record whether both perturbations match the base structural metadata and
  the realized acceptance path.
  \item If either side changes branch, mark the row branch-invalid and do not
  use it as derivative evidence.
  \item If the row is branch-valid, compute \(D_i(h)\) and the relative error
  \(e_i(h)\).
  \item Report whether two consecutive branch-valid rows meet the declared
  tolerance rule.
\end{enumerate}
\end{minipage}}

\section{Validation, Diagnostics, and What Counts as Evidence}
\label{sec:validation}

FixedSGQF should be validated as an approximate deterministic likelihood lane,
not as exact posterior propagation.

\paragraph{Implementation-correctness evidence.}
The minimum implementation audit should include:
\begin{enumerate}[leftmargin=0.3in]
  \item cloud-construction checks, including duplicate-node accumulation, weight
  totals, negative-weight counts, and exact reproduction of the toy cloud in
  Section~\ref{subsec:toy-cloud};
  \item linear-Gaussian recovery, where the Gaussian projection recursion should
  match the Kalman filter up to quadrature and floating-point tolerance;
  \item the worked one-step numerical oracle in
  Section~\ref{sec:worked-example};
  \item same-scalar finite differences on representative parameter points;
  \item signed-weight stability diagnostics for the minimum eigenvalues of
  symmetrized \(P_t^-\), \(S_t\), and \(P_t\), plus any branch veto.
\end{enumerate}

\paragraph{Approximate scientific-usefulness evidence.}
Separate from implementation correctness, one may also report:
\begin{enumerate}[leftmargin=0.3in]
  \item accuracy ladders across sparse-grid levels,
  \item comparisons against dense quadrature in small dimension,
  \item comparisons against linear-Gaussian truth when available,
  \item comparisons against particle or tensor-density references in regimes
  where richer posterior structure matters,
  \item runtime, point count, and memory scaling.
\end{enumerate}
These latter diagnostics support judgments about usefulness of the chosen
approximation; they do not by themselves certify correct implementation of the
saved scalar.

\section{Neighboring Methods and Why FixedSGQF Is the Selected Lane}
\label{sec:comparison}

Nearby filters can be compared by the carried object, the moment engine, and
whether they admit the deterministic fixed-scalar derivative contract used here.

\begin{center}
\small
\begin{tabular}{p{0.15\linewidth}p{0.21\linewidth}p{0.20\linewidth}p{0.17\linewidth}p{0.19\linewidth}}
\toprule
method & carried object & moment engine & same-scalar gradient lane & main limitation \\
\midrule
EKF & one Gaussian & local linearization & yes, but for the linearized moment engine & weak nonlinear moment fidelity \\
UKF/CKF & one Gaussian & fixed low-order deterministic points & yes & lower-order point design \cite{julier1996general,vandermerwe2004sigma,arasaratnam2009cubature} \\
Dense GHQF & one Gaussian & dense tensor-product Gaussian quadrature & yes & point count grows exponentially in dimension \\
FixedSGQF & one Gaussian & fixed sparse-grid quadrature & yes & still one Gaussian; branch-sensitive signed-weight moments \cite{jia2012sparsegrid} \\
Adaptive sparse grid & one Gaussian & parameter-dependent active grid & only piecewise, unless frozen first & target changes if the live grid changes \cite{singh2018adaptive} \\
Particle / TT lanes & richer non-Gaussian object & samples or density representation & different target class & heavier approximation object; different contract \cite{zhao2024ttsequential} \\
\bottomrule
\end{tabular}
\end{center}

The selection argument of this note is local rather than universal.  FixedSGQF
is selected here because it jointly provides:
\begin{enumerate}[leftmargin=0.3in]
  \item a declared deterministic approximate likelihood scalar,
  \item an analytical derivative of that same scalar on a fixed smooth branch,
  \item a moment engine less explosive than dense tensor-product Gaussian
  quadrature in the intended high-dimensional regime.
\end{enumerate}
This is not a theorem that FixedSGQF is the best nonlinear filter overall.  It
is the selection statement for the deterministic Gaussian-surrogate lane treated
in this note.

\section{Conclusion}
\label{sec:conclusion}

This note has defined an implementation-ready mathematical companion for fixed
sparse-grid quadrature filtering.  The document states the exact target and the
narrower carried object, constructs the fixed sparse-grid cloud in standardized
coordinates, gives the one-step value recursion, fixes the saved-branch
same-scalar contract, derives the branch-conditioned analytical gradient, and
states the finite-difference rule needed to check that value and derivative refer
to the same deterministic scalar.

The note has also preserved the central non-claim.  FixedSGQF is not exact
nonlinear filtering, and it does not preserve general non-Gaussian posterior
shape.  It is a deterministic Gaussian-surrogate lane.  Its approval claim is
therefore bounded: the claim is not universal optimality, but that this lane is
coherent, mathematically explicit, auditable by implementation engineers, and
sufficiently explicit for mathematically trained junior workers to implement at
the algorithm level.

\appendix

\section{Implementation-Oriented Notes for a CS-Facing Appendix}
\label{sec:appendix-cs}

This appendix records software-facing detail that is helpful for implementation
engineers but is not mathematically indispensable to the main note.

\subsection{Suggested record layout}

A practical implementation may distinguish at least four record types:
\begin{enumerate}[leftmargin=0.3in]
  \item a cloud record containing standardized nodes, accumulated weights,
  negative-weight diagnostics, and deterministic ordering metadata;
  \item a one-step value cache containing placed points, transition images,
  observation images, and moment objects reused by the gradient path;
  \item a branch record containing all fixed structural policies used to define
  the saved scalar;
  \item a failure record of the form
  \((\text{time},\text{stage},\text{reason},\text{diagnostic values})\).
\end{enumerate}

\subsection{Deterministic storage and ordering}

Because node ordering is part of the saved scalar contract, implementations
should store the merged cloud in a deterministic sorted order after duplicate
accumulation.  Hash-based storage may be used during construction, but the final
serialized object should record the deterministic order actually used by value,
gradient, and finite-difference routines.

\subsection{Suggested routine surfaces}

A software-facing implementation may expose routines analogous to:
\begin{enumerate}[leftmargin=0.3in]
  \item \texttt{build\_fixed\_sparse\_grid\_cloud},
  \item \texttt{fixed\_sgqf\_value},
  \item \texttt{fixed\_sgqf\_gradient},
  \item \texttt{same\_scalar\_fd\_check}.
\end{enumerate}
The exact API is an engineering choice, but these routine boundaries should
remain consistent with the mathematical contract of the main note.

\bibliographystyle{plainnat}
\bibliography{../references}

\end{document}
